
\section{Other Attempts at Modeling PubChem AID 1851}
PubChem AID 1851 has been the basis for several attempts to advance \textit{in silico} screening development. Several research groups have developed large-scale single target models for individual isozymes~\cite{Vasanthanathan2009,Sridhar2012,Sun2012,Novotarskyi2010}. Another group pursued all five isozymes but used a proteochemometric method that also takes into account each isozyme's protein sequence~\cite{Lapins2013}. Finally, Cheng et al. built single-target QSAR models for all five CYP isozymes, the work most comparable to this thesis~\cite{Cheng2011}. All of these models exceeded the 0.6 accuracy threshold used in this thesis.

Cheng et al. modeled the same five isozymes but used a combined classifier on a larger data set of more than 24,700 unique compounds. These compounds also came from PubChem AID 1851 and included the 17,143 used in this study, plus additional compounds from PubChem AID 410. The group used an ensemble of different independent machine learning classifiers including support vector machine, C 4.5 decision tree, \textit{k}-nearest neighbors, and naive Bayes. The final ensemble was brought together using a back-propagation artificial neural network (BP-ANN). Those models were also internally validated by 5-fold cross-validation and then evaluated with a test set composed of about 9000 diverse unique compounds previously unseen by the classifier. Test set validation measures ranged from 0.764 to 0.886, somewhat higher than the test accuracies reported here (0.674--0.804), with a larger data set and an ensemble method. Cheng et al. claim these classification models are applicable for virtual screening for inhibitors of the five major CYP isozymes or can be used as simple filters of potential chemicals in drug discovery~\cite{Cheng2011}.


\section{Sources of Error}
This assay required recombinant sources of CYP enzymes because the probes are not sufficiently P450 isozyme-selective to be used with those derived from human liver microsomes. Another potential source for false positives in these assays can be compounds that interfere with light generation directly, such as compounds that interfere with luciferase enzymatic activity~\cite{Zlokarnik2005}.

The results of the large-scale screening of PubChem AID 1851 against five CYP isozymes found that most compounds in a typical chemical library inhibited several isozymes, and only a small fraction inhibited none~\cite{Veith2009}.

The uniformly lower accuracy on non-inhibitors than on inhibitors may reflect the composition of the non-inhibitor class, which contains both true non-inhibitors (score 0) and inconclusive compounds (scores 1--39). A binary labelling may discriminate less well than one with more classes. Future models could test this by adding a separate inconclusive class or changing the Activity Score threshold.


\section{Machine Learning for Pharmaceutical Sciences}
\begin{quote}
Developing successful machine learning applications still requires a substantial amount of “black art” that is hard to find in textbooks~\cite{Domingos2012}.
\end{quote}

Unlike the data in this study, raw data are often not amenable to learning, but features constructed from them can be. Features are the most important factor in a good model. If there are many independent features that correlate well with the class, learning is easy. However, if the class is a very complex function of the features, training a good model may take too long or be impossible~\cite{Domingos2012}. Most of the effort in a machine learning project goes into feature engineering, an area where domain expertise and experience with chemical library design can mean the difference between success and failure.

Machine learning is not a one-shot process of building a data set and running a learner, but rather an iterative process of running the learner, analyzing results, modifying the data and/or learner, and repeating. Learning is often the quickest part of this, but only because of all the effort that has gone into sharing and codifying prior knowledge. Feature engineering is considered more difficult because it is domain-specific, while learners can by and large be general-purpose~\cite{Domingos2012}.

Domain expertise in feature engineering has no substitute. The alternative, running a learner on a very large number of features to find useful combinations, may be too time-consuming or cause overfitting~\cite{Eklund2014}.

As a rule, it pays to try the simplest learners first (e.g. naive Bayes before logistic regression, \textit{k}-nearest neighbors before support vector machines). More sophisticated learners may be seductive, but they are usually harder to use, because there are more tuning parameters and their internals are more opaque~\cite{Domingos2012}.

One way to divide learning algorithms is as follows: those whose representation has a fixed size, like linear classifiers, and those whose representation can grow with the data, like decision trees. Fixed-size learners can only take advantage of so much data. Variable-sized learners can in principle learn any function given sufficient data, but in practice may not, because of limitations of the algorithm or computational cost. Also, because of the curse of dimensionality, no existing amount of data may be enough~\cite{Domingos2012}. Each type of learner comes with its own assumptions and, consistent with the no free lunch theorem, none is best in all situations, although some are clearly better than most~\cite{Hand2006,Delgado2014}.

Perhaps because different learners see the data from different angles, combining them into an ensemble often improves accuracy, and ensembles are now standard in machine learning. Techniques such as bagging, boosting and stacking resample the training data or reweight previously misclassified instances to reduce variance while minimizing bias. The simplest ensemble is a majority vote: each instance is assigned the class predicted by most of the models.


Because ensembles frequently outperform their constituent models, a logical next step is to combine the models developed here.

